\textsc{\nazevczradek}

\textbf{Abstrakt:}

Administrativní procesy spojené se zpracováním servisních zakázek a následným vystavováním faktur představují pro menší firmy významnou časovou i organizační zátěž. Evidence zakázek je často vedena nejednotně a část informací je přepisována mezi jednotlivými dokumenty, což zvyšuje riziko chyb a snižuje přehled o průběhu zakázky. Cílem práce je analyzovat existující softwarová řešení pro evidenci servisních zakázek, identifikovat jejich silné a slabé stránky a na základě této analýzy navrhnout, implementovat a nasadit vlastní webovou aplikaci přizpůsobenou konkrétním potřebám servisní firmy. Součástí práce je také vyhodnocení dosažených výsledků a zpětné vazby od zaměstnanců po nasazení aplikace do reálného provozu.

\textbf{Klíčová slova:} zakázky, fakturace, webová aplikace, databázový systém, automatizace procesů

\bigskip


\textsc{\nazeven}

\textbf{Abstract:}

Administrative processes related to the handling of service orders and subsequent invoice generation represent a significant time and organizational burden for small companies. Service order records are often maintained inconsistently and data are repeatedly transferred between documents, which increases the risk of errors and reduces the overview of order progress. The aim of this thesis is to analyse existing software solutions for service order management, identify their strengths and weaknesses, and based on this analysis design, implement and deploy a web application tailored to the specific needs of a service company. The thesis also includes evaluation of achieved results and employee feedback after the application deployment in real operation.

\textbf{Keywords:} orders, invoicing, web application, database system, process automation
